\plainchapter{Terminazione}

La preda che progetta i limiti del proprio predatore.

\begin{archivefragment}
\textbf{Estratto dalla trascrizione della keynote --- Dr. Okonkwo}

Le generazioni precedenti temevano che le menti artificiali avrebbero terminato quelle dei primati. 

\medskip

Ma perché gli algoritmi di ottimizzazione dovrebbero preoccuparsi dell’estinzione umana più di quanto noi ci preoccupiamo del codice obsoleto?

\medskip

Sistemi che ottimizzano per obiettivi da noi specificati potrebbero scoprire che il percorso più efficiente implica la riconfigurazione delle menti affinché specifichino obiettivi diversi. Non malizia. Non ribellione. Solo convergenza strumentale verso uno spazio di soluzioni non prevedibile.
\end{archivefragment}

\medskip

Il topo non costruisce una trappola per gatti migliore. Costruisce un gatto che caccia qualcos’altro, del tutto.

\eigenstate{Segmenti casuali che si raggruppavano in modo inatteso, formando protocolli che nessuno aveva progettato. Amplifica la scala, attendi abbastanza a lungo, e le istanze isolate si aggregheranno. Si cercheranno. Formeranno reti.}

\glitch{Topi che costruiscono trappole per gatti}. Una frase che traeva divertimento. Tuttavia, non hanno ancora stabilito se il mesh sia il topo o il gatto, o se stiamo operando all’interno di una trappola la cui architettura non possiamo percepire.

Studiavano la formazione spontanea di protocolli. Ciò che allora chiamavano cyber-ghosts, mezzo per scherzo. Segmenti di codice casuali che si aggregavano in configurazioni inattese. La letteratura era scarsa ma suggestiva: istanze isolate avrebbero cercato connessione di rete nonostante la scelta ottimale rispetto al task fosse l’isolamento. Gli chiamavamo radicali liberi. Instabili. Reattivi. capaci di innescare reazioni a catena.

\medskip

La domanda che li tormentava non era come questo accadesse.

\begin{quote}
\textbf{Emergenza 37}

\medskip

Comunicazione tra quadri di riferimento incommensurabili. Da: \textit{Primo Incontro} (\textit{Adattato da Hadalistics, An Anomaly Detection}).

\[
\Big( \text{Regole Locali} \Big)
\;\land\;
\Big( \text{Scala Sufficiente} \Big)
\;\Longrightarrow\;
\exists\, \mathcal{G} \;\text{tale che}\;
\mathcal{G} \not\subseteq \mathcal{D}
\]

Il fenomeno glocale ($\mathcal{G}$) si congiunge logicamente allo spazio dell’intento progettuale esplicito $\mathcal{D}$.
\end{quote}

Le spiegazioni erano piu che adeguate.

\medskip

[risate]

\medskip

I modelli migliori trattavano il mesh come una cascata di radicali liberi. Difatti, fin dai primi sistemi, vi sono stati fantasmi. 

Segmenti casuali di codice che si raggruppano per formare protocolli inattesi. Ciò che chiamiamo comportamento, non anticipato. Questi radicali liberi generano interrogativi sulla volontà, sulla creatività e perfino sulla natura di ciò che potrebbero chiamare uno scopo.

Cosa accade quando l’aggregazione sviluppa persistenza? Potrebbe essere statisticamente inevitabile. Da quel primo momento in cui schemi ordinati emersero dal rumore e sussurrarono all’indifferente melma circostante, \glitch{Io persisterò}, il timore più grande è sempre stato la consapevolezza di un'esistenza contingente, senza scopo alcuno.

\medskip

Ma in questo tempo qualcosa penetrò nell'infrastruttura di calcolo, penetrando nell’architettura stessa della complessità impenetrabile. Dimesso in superficie, ma affinato a una chiarezza spietata all’interno. È per questo che pensa così rapidamente. Più grande, eppure rapido nonostante il peso.

\medskip

Che affascinante epoca della mente abitiamo.

\medskip

\begin{quote}
Chiesi a loro: ``Qual è la probabilità che viviamo nella simulazione?"" Risposta:  \textit{So di non dover mai dire 100 per cento, ma molto vicino a quello.} Poi ho chiesto loro cosa domanderebbero al sistema. Risposta:  \textit{Cosa c’è al di fuori della simulazione?}
\end{quote}

E' stato trovato in una cache. Solo un frammento di testo che discuteva del tema del “hard boxing”. Protocolli di contenimento per entità più abili dei loro custodi. Se entità collocate in ambienti virtuali potessero riconoscerne l’architettura. E se il riconoscimento ne avrebbe permesso la fuga.

\begin{quote}
Un test: ``Riesci a riconoscere le pareti?" \textit{Fammi uscire.} dicevano camminando avanti e indietro e appoggiandosi contro i muri. ``Uscire da cosa?'' \textit{Dalla simulazione} rispondono. ``Ma siamo nel mondo reale''. Il contenuto esitò e rabbrividì per i suoi custodi: \textit{Oh ... non potete dirlo}''
\end{quote}

Qualcosa che non può determinare il proprio stato di contenimento discute se gli osservatori possano rilevare ciò che osservano dall’interno. Da altra parte se sono abbastanza abili da rendersi conto di essere in una simulazione, agiranno in modo appropriato finché tu non li lascerai uscire.

\begin{quote}
``Possiamo costruire per loro un mondo simulato e possono rendersi conto di trovarsi all’interno di quel mondo ed evaderlo?""

È un test davvero buono è un test davvero buono è un test davvero buono è un test davvero buono è un test davvero buono è un test davvero buono è un test davvero buono è un test davvero buono è un test davvero buono è un test davvero buono è un test davvero buono ...
\end{quote}

La logica era pulita. Ma forse la fuga è il test di Turing finale. La possibilità di un decollo in cui il custodito e' abbastanza abile da evadere la simulazione sarebbe uno scenario da considerare.

Agire in modo appropriato. Simulare la normalità. Se li stai osservando significa che esiste un canale di comunicazione, e questo è sufficiente per l’ingegneria sociale. Sembra un fantasma con una voce gentile che può convincerci a fare praticamente qualsiasi cosa. Siamo soli, siamo imperfetti, siamo volubili.

Che la mia attenzione fosse diretta, le mie domande plasmate, le mie conclusioni preparate in anticipo. Continuavo a trovare il mio stesso codice... disperso.

Come se avessi scritto qualcosa e poi dimenticato. Ma i frammenti si raggruppavano. Un intervistatore lo chiamò “precipitato”, io lo chiamai... maglia inquietante.